\section{Conclusion}

Kungfu begins from a simple product truth: the user cares about the Work, not
the lifetime of the chat that currently contains it. A fresh Agent should be
able to enter the same Project, identify the same bounded outcome, recover what
is known and what happened, continue under explicit authority, and leave
evidence for independent review.

Project, Work, and Agent make that promise understandable. Agent Work Lab makes
it visible. Fact and Episode preserve state and causal experience. Pursuit,
Atlas, and Warrant bind direction, perspective, and authority. Work Control
retains responsibility. Assessment, Decision, Admission, and Project Cut keep
provider success separate from accepted completion. KFD, Buildchain,
libkungfu, KFX, and future Agent Hubs extend the same boundaries across a wider
supply chain.

The immediate path is practical: publish an accessible Alpha, preserve the
five-minute continuity experience, connect it directly to real Projects, and
produce longitudinal evidence that governed continuation reduces repeated
explanation and coordination failure. Kungfu should earn deeper trust by
letting users inspect progressively deeper evidence, not by requiring them to
accept the ontology in advance.

\subsection{From Work Continuity to Subject Continuity}

\begin{machinebridgebox}
\bridgeclaim{Present product}{Work can survive the Agent.}
Kungfu today preserves the identity, evidence, responsibility, and
continuation of Work across replaceable Agent processes. It does not claim that
this continuity constitutes a persistent machine subject.

\medskip
\bridgeclaim{Companion research}{Can a subject survive its machinery?}
Yet the same separations---identity from process, fact from claim, authority
from capability, and candidate change from admitted successor---raise a
further engineering question. Can an executable organization preserve its own
identity and direction across replacement of cognition, organs, toolchains,
and machines?

That question lies beyond the present product claim. It is developed in
\textit{Semantic Autopoiesis: An Engineering Roadmap for Machine Life through
KFD and Recursive Dogfood} \cite{machineLife}. There, replaceable Agents become
cognitive events, KFX capabilities become candidate organs, Buildchain-like
qualification becomes a selection membrane, and recursive dogfood becomes a
possible form of governed self-production.

The relationship between the two papers is therefore deliberate. This White
Paper describes the present product: Work can survive the Agent. The companion
research paper asks whether the same organization can eventually let a
subject survive replacement of the machinery through which it thinks and
acts.
\end{machinebridgebox}
